\documentclass[11pt,a4paper,fontset=fandol]{ctexart}
\usepackage[a4paper,margin=1.78cm,headheight=15pt]{geometry}
\usepackage{amsmath,amssymb,mathtools,bm}
\usepackage{booktabs,longtable,tabularx,array}
\usepackage{enumitem}
\usepackage[most]{tcolorbox}
\usepackage{tikz}
\usetikzlibrary{arrows.meta,positioning,fit,calc}
\usepackage{xcolor}
\usepackage{fancyhdr}
\usepackage{hyperref}
\usepackage{bookmark}
\usepackage{microtype}

\newcolumntype{Y}{>{\raggedright\arraybackslash}X}
\newcolumntype{L}[1]{>{\raggedright\arraybackslash}p{#1}}
\setlength{\parindent}{2em}
\setlength{\parskip}{0.34em}
\setlength{\emergencystretch}{3em}
\renewcommand{\arraystretch}{1.22}
\setlength{\tabcolsep}{3pt}
\setlist[itemize]{itemsep=0.18em,topsep=0.35em,leftmargin=2em}
\setlist[enumerate]{itemsep=0.18em,topsep=0.35em,leftmargin=2em}

\definecolor{deep}{RGB}{42,58,73}
\definecolor{soft}{RGB}{245,247,249}
\definecolor{accent}{RGB}{104,76,55}
\definecolor{greenish}{RGB}{57,92,76}
\definecolor{warn}{RGB}{136,68,63}
\definecolor{purple}{RGB}{87,72,116}
\hypersetup{colorlinks=true,linkcolor=deep,urlcolor=accent,pdftitle={一元累积：原函数、定积分与反常积分}}
\pagestyle{fancy}
\fancyhf{}
\lhead{\small\color{deep}\textbf{数学一 · 高等数学 Topic 05}}
\rhead{\small\color{deep}原函数 · 定积分 · 反常积分}
\cfoot{\small\thepage}
\renewcommand{\headrulewidth}{0.4pt}
\newtcolorbox{corebox}[1]{breakable,colback=soft,colframe=deep,title=\textbf{#1},boxrule=0.8pt,arc=2mm}
\newtcolorbox{methodbox}[1]{breakable,colback=white,colframe=greenish,title=\textbf{#1},boxrule=0.7pt,arc=1.5mm}
\newtcolorbox{warnbox}[1]{breakable,colback=white,colframe=warn,title=\textbf{#1},boxrule=0.7pt,arc=1.5mm}
\newtcolorbox{examplebox}[1]{breakable,colback=white,colframe=purple,title=\textbf{#1},boxrule=0.7pt,arc=1.5mm}
\newcommand{\kw}[1]{\textbf{\color{deep}#1}}

\begin{document}
\begin{center}
{\LARGE\textbf{一元累积：原函数、定积分与反常积分}}\\[0.4em]
{\large 从“局部贡献”到“有限总量与无限尾部”}
\end{center}
\vspace{0.5em}

\begin{quote}
本册只追问一个根本问题：\textbf{把无穷多个局部贡献累加成一个可解释的整体时，究竟要先固定对象、区间与正则性，再用什么方式计算和取极限？} 原函数是反向找变化率的工具，定积分是有限区间的有向累积，反常积分则是带奇点或无穷端点的截断极限；三者相关，但不能用一个公式抹平它们的资格边界。
\end{quote}

\begin{center}\rule{0.5\linewidth}{0.5pt}\end{center}

\setcounter{section}{-1}
\section{本册定位}

本册是高等数学 Subject Atlas 下的 Topic05。Topic03 提供局部导数，Topic04 提供区间形状；本册把局部贡献组织为有限累积、物理/几何总量和可审计的无限尾部。

\subsection{Own}

本册独立负责：
\begin{itemize}
  \item 原函数族、不定积分的对象边界与反导方法的选择；
  \item 定积分的线性、区间可加、保号、估值、对称和换元结构；
  \item 定积分计算的预处理顺序、积分不等式与平均高度视角；
  \item 微元法在面积、弧长、旋转体、功、压力与形心中的统一翻译；
  \item 反常积分的奇点识别、区间拆分、截断计算与逐段取极限；
  \item 最终保号与持续振荡两类判敛机制，以及绝对/条件收敛边界；
  \item “先有限、后极限”的累积协议与一元积分题的失败路径。
\end{itemize}

\subsection{Uses / Stop Boundary}

本册调用 Topic01 的定义域和函数对象，调用 Topic02 的极限/等价/夹逼语言，调用 Topic03 的导数与有限 Taylor，调用 Topic04 的单调性、凹凸和中值定理。以下内容只保留接口：
\begin{itemize}
  \item 原函数、Riemann 可积、变上限积分可导性的细致正则性分类归 H-B03；本册只使用其结论接口；
  \item 反常积分与数项级数的共同尾部结构、积分判别和连续/离散翻译归 H-B04；
  \item 有限 Taylor 的局部余项归 Topic03，Taylor 级数的无限表示与收敛域归 Topic11/H-B05；
  \item 多重积分、曲线曲面积分和完整微分方程解法不在本册展开；
  \item 题目出现何种信号、先做哪一种预处理、何时停止计算，进入《高等数学做题规则》。
\end{itemize}

本册解释累积机制为什么成立；Rules 维护考场上的触发器、第一动作、检查点和失效分支。公式表只保留能重建结构的母式，不把整本归档笔记复制为第二个知识 Owner。

\section{根本问题：局部贡献怎样成为整体？}

设一个量在位置 $x$ 附近的微小贡献近似为 $q(x)\,dx$。积分不是“把符号 $\int$ 写上去”，而是先切成有限个小区间，再让最大网格趋于零，检查所有取样方式的总和是否稳定：
\[
S(P,\boldsymbol\xi)=\sum_i q(\xi_i)\Delta x_i,
\qquad \lVert P\rVert=\max_i\Delta x_i\to0.
\]
若总和趋于同一个数，才得到
\[
\int_a^b q(x)\,dx.
\]

因此本册的统一对象不是“积分技巧”，而是四层审计：
\[
\boxed{\text{对象} \to \text{合法区间} \to \text{有限累积} \to \text{极限与解释}}.
\]

\begin{corebox}{累积的母模型}
任何一元积分题都可重建为
\[
\boxed{\begin{gathered}
\text{Local Contribution}
\to\text{Domain / Regularity}
\to\text{Partition / Transform}\\
\text{Finite Accumulation}
\to\text{Boundary / Limit}
\to\text{Verification}
\end{gathered}}
\]
局部贡献决定被积函数；定义域与正则性决定能否积；分割、对称和换元决定如何降低结构复杂度；先在正常区间上完成有限计算，再处理端点或无穷的极限；最后用符号、单位、数量级和反向求导独立核验。
\end{corebox}

\begin{center}
\begin{tikzpicture}[
  node distance=0.55cm and 0.55cm,
  box/.style={draw=deep,rounded corners=1.5mm,minimum width=3.05cm,minimum height=1.0cm,align=center,fill=soft,font=\small},
  arr/.style={-{Latex[length=2mm]},thick,draw=deep}
]
\node[box] (a) {局部贡献\\$q(x)\,dx$};
\node[box,right=of a] (b) {对象与区间\\可积/连续?};
\node[box,right=of b] (c) {分割与变换\\降结构复杂度};
\node[box,below=of c] (d) {正常积分\\有限累加};
\node[box,left=of d] (e) {端点/无穷\\逐个取极限};
\node[box,left=of e] (f) {解释与核验\\符号/单位/反导};
\draw[arr] (a)--(b); \draw[arr] (b)--(c); \draw[arr] (c)--(d);
\draw[arr] (d)--(e); \draw[arr] (e)--(f);
\end{tikzpicture}
\end{center}

\section{三个对象不能混成一个“积分”}

\subsection{原函数是反向变化率}

若在同一个连通区间 $I$ 上 $F'(x)=f(x)$，则 $F$ 是 $f$ 的一个原函数，所有原函数组成
\[
\int f(x)\,dx=F(x)+C.
\]
若 $F,G$ 在同一\textbf{区间}上都是原函数，则 $(F-G)'=0$，所以 $F-G=C$。区间若被断开，各连通分支可以有不同常数；这正是 $\int dx/x=\ln|x|+C$ 不能跨过 $0$ 统一成一个单常数表达的原因。

“存在原函数”与“能写出初等原函数”不是同一问题。例如 $e^{-x^2}$ 在连续区间上有原函数，即使该原函数通常不能用初等函数表达。计算题可以改用定积分定义该原函数，不能因此宣称积分不存在。

\subsection{Riemann 积分是区间累积}

在 $[a,b]$ 上，Riemann 可积要求所有取样方式的和在网格趋零时趋于同一数。连续函数自动满足这一条件；有界、分段连续等是常用可积来源。点值在有限个点上的修改不改变定积分，但会改变“逐点导数”等局部问题，故不能把积分等价误用为函数对象等价。

\subsection{变上限函数是新对象}

固定 $a$，定义
\[
A(x)=\int_a^x f(t)\,dt.
\]
这里 $t$ 是累积变量，$x$ 是新函数的上限。即使 $f$ 可积，也要另行检查 $A$ 的连续性和 $A'(x)=f(x)$ 是否成立。连续 $f$ 给出标准的基本定理接口；一般可积但不连续的情形由 H-B03 的平均值商判据接管。

\begin{warnbox}{“积分与求导互逆”必须写对象和条件}
\[
\frac{d}{dx}\int_a^x f(t)\,dt=f(x)
\]
不是无条件的符号消去。要先问：被积函数在哪个区间讨论、是否在该点连续、上下限是否依赖 $x$、积分变量与函数变量是否混淆。反方向的 Newton\textendash{}Leibniz 公式也要求已经有合资格的原函数或基本定理条件。
\end{warnbox}

\section{有限累积的代数结构与预处理}

\subsection{四条基础性质}

在相应区间可积时，定积分满足
\[
\int_a^b(\alpha f+\beta g)=\alpha\int_a^b{f}+\beta\int_a^b{g},
\qquad
\int_a^b{f}=\int_a^c{f}+\int_c^b{f},
\]
\[
\int_a^b{f}\ge0\quad(f\ge0),
\qquad
\left|\int_a^b{f}\right|\le\int_a^b{|f|}.
\]
若 $m\le f(x)\le M$，则
\[
m(b-a)\le\int_a^b{f(x)}\,dx\le M(b-a).
\]
这说明积分是一个保序的线性累积器；比较积分时优先比较平均高度，而不是盲目求出两个复杂原函数。

\subsection{计算前的结构扫描}

遇到定积分，按如下顺序扫描：
\begin{enumerate}
  \item 区间是否对称、可翻转或由若干周期组成；
  \item 被积函数是否含绝对值、取整、最值或分段点，是否必须先拆区间；
  \item 是否能通过 $x\mapsto a+b-x$、平移或三角换元配对；
  \item 能否直接作差、作商或用保序性/估值完成目标；
  \item 仍需计算时，选择凑微分、换元、分部积分或递推；
  \item 最后才调用原函数，并反向求导或做数量级检查。
\end{enumerate}

\subsection{对称、周期与区间翻转}

对称区间上奇函数积分为零、偶函数积分为半区间两倍。一般区间满足
\[
\int_a^b{f(x)}\,dx=\int_a^b{f(a+b-x)}\,dx,
\]
故可平均为
\[
\int_a^b{f(x)}\,dx
=\frac12\int_a^b\bigl[f(x)+f(a+b-x)\bigr]dx.
\]
若 $f$ 周期为 $T$，每个长度为 $T$ 的区间积分相同，整数个周期可先压缩为一个周期。区间翻转只是变量替换，不会自动消除换元后的微分因子。

\subsection{定积分换元与分部积分}

单调可导换元 $x=\phi(t)$ 时，必须同时变换上下限与微分：
\[
\int_a^b{f(x)}\,dx=\int_\alpha^\beta{f(\phi(t))\phi'(t)}\,dt,
\qquad \phi(\alpha)=a,\ \phi(\beta)=b.
\]
分部积分保持边界项：
\[
\int_a^b u\,dv=\left.uv\right|_a^b-\int_a^b{v}\,du.
\]
边界项是结构的一部分，不能像不定积分那样暂时省略；在反常积分中还必须先在截断区间写出，再分别取极限。

\section{反导工具箱：方法由结构决定}

不定积分的核心不是背更长的公式，而是识别导数结构：
\[
\begin{tikzpicture}[>=Stealth]
  \tikzset{
    flowstage/.style={draw=deep, fill=white, rounded corners=1.5mm, minimum height=0.72cm, align=center, font=\small\sffamily\bfseries\color{deep}, line width=0.6pt, inner xsep=6pt, inner ysep=3.5pt},
    flowarr/.style={->, >=Stealth, line width=0.75pt, draw=accent}
  }
  \node[flowstage] (n1) at (0,0) {1. 先预处理};
  \node[flowstage, right=0.4cm of n1] (n2) {2. 凑微分};
  \node[flowstage, right=0.4cm of n2] (n3) {3. 换元};
  \node[flowstage, below=0.6cm of n3] (n4) {4. 分部};
  \node[flowstage, left=0.4cm of n4] (n5) {5. 递推/部分分式};
  \node[flowstage, left=0.4cm of n5] (n6) {6. 反向求导};

  \draw[flowarr] (n1) -- (n2);
  \draw[flowarr] (n2) -- (n3);
  \draw[flowarr] (n3) -- (n4);
  \draw[flowarr] (n4) -- (n5);
  \draw[flowarr] (n5) -- (n6);
\end{tikzpicture}
\]

\subsection{凑微分与换元}

若被积式包含 $u'(x)\Phi(u(x))$，令 $u=u(x)$，把问题压缩为 $\int\Phi(u)\,du$。常见母式包括
\[
\int\frac{u'}u\,dx=\ln|u|+C,
\qquad
\int u^\mu u'\,dx=\frac{u^{\mu+1}}{\mu+1}+C\quad(\mu\ne-1).
\]
换元前必须确认 $u$ 在所讨论区间有定义；根式、对数和分母的合法性先于形式积分。

\subsection{分部积分的选择}

分部积分把“容易求原函数的部分”与“求导后更简单的部分”交换：
\[
\int u\,dv=uv-\int v\,du.
\]
当一个因子求导后降阶（多项式），另一个因子可直接反导（指数、三角）时，优先使用。若循环回原积分，整理边界项后再解一次代数方程；若递推未降低复杂度，应停止继续分部。

\subsection{有理、根式与三角结构}

有理函数先作多项式除法，再对真分式部分做部分分式；根式题先找能把根式变成有理函数的替换；三角幂次按奇偶性保留一个可凑微分因子或使用降幂。每次选择都由“变换后复杂度是否下降”决定，而不是按题型名称机械套用。

\section{比较积分：从有向面积到平均高度}

比较
\[
I=\int_a^b{f(x)}\,dx,\qquad J=\int_c^d{g(x)}\,dx
\]
时，第一动作是统一区间。换元后应比较完整的密度 $g(\phi(t))\phi'(t)$，不能只看函数值。若区间一致，优先作差
\[
I-J=\int_a^b(f-g)\,dx,
\]
并用符号、单调性、凸性、局部放缩或配对判断。若 $f\ge g$ 且连续并非处处相等，则严格不等式来自一段正面积，而不是来自某个孤立点。

\subsection{作差、作商与局部放缩}

作差适合符号直接可见的结构；作商适合同号正函数的相对增长；局部放缩适合只需一个统一上界或下界的估值。若出现 $e^x\ge1+x$、$\ln(1+x)\le x$、$\sin x\le x$ 等母不等式，先检查区间和方向，再积分。若一阶放缩发生主部相消，调用 Topic03 的有限 Taylor 接口，保留相消后的第一项与余项阶数。

\subsection{积分中值与凸性}

若 $f$ 在 $[a,b]$ 连续，则存在 $\xi\in[a,b]$ 使
\[
\int_a^b{f(x)}\,dx=f(\xi)(b-a).
\]
它说明积分等于“某个平均高度 $\times$ 区间长度”，但只产生存在性点，不能把 $\xi$ 固定为中点。若 $f$ 凸，则梯形平均与中点值之间的方向由切线/割线决定；需要更强的权重不等式时再调用相应 Bridge，不把未验证的 Jensen 或 Chebyshev 结论当成通用规则。

\section{微元建模：先定义贡献，再累加}

几何/物理应用的稳定流程是：
\[
\begin{tikzpicture}[>=Stealth]
  \tikzset{
    flowstage/.style={draw=deep, fill=white, rounded corners=1.5mm, minimum height=0.72cm, align=center, font=\small\sffamily\bfseries\color{deep}, line width=0.6pt, inner xsep=6pt, inner ysep=3.5pt},
    flowarr/.style={->, >=Stealth, line width=0.75pt, draw=accent}
  }
  \node[flowstage] (n1) at (0,0) {1. 对象与坐标};
  \node[flowstage, right=0.4cm of n1] (n2) {2. 切片/微元};
  \node[flowstage, right=0.4cm of n2] (n3) {3. 局部贡献};
  \node[flowstage, below=0.6cm of n3] (n4) {4. 积分区间};
  \node[flowstage, left=0.4cm of n4] (n5) {5. 单位与边界核验};

  \draw[flowarr] (n1) -- (n2);
  \draw[flowarr] (n2) -- (n3);
  \draw[flowarr] (n3) -- (n4);
  \draw[flowarr] (n4) -- (n5);
\end{tikzpicture}
\]

\subsection{面积、弧长与旋转体}

竖直切片给 $dA=[\text{上函数}-\text{下函数}]dx$；水平切片则用 $dy$。弧长来自微小线段
\[
ds=\sqrt{dx^2+dy^2}=\sqrt{1+{(y')}^2}\,dx.
\]
绕坐标轴旋转时，圆盘/垫圈法累加截面积，圆柱壳法累加周长乘高：
\[
V=\pi\int(R^2-r^2)\,dx,
\qquad
V=2\pi\int(\text{半径})(\text{高})\,dx.
\]
选择切片方向的标准是：让半径、高度和积分区间由同一个自变量直接表达，避免多余反函数分支。

\subsection{功、压力与形心}

变力做功是
\[
dW=F(x)\,dx,\qquad W=\int_a^b{F(x)}\,dx.
\]
抽水做功的微元是“微小重量 $\times$ 提升距离”；静水压力的微元是“压强 $\rho gh\times$ 条带面积”。形心/质心的统一思想是先累加质量与一阶矩，再取比值：
\[
\bar x=\frac{\int x\,dm}{\int dm},
\qquad
\bar y=\frac{\int y\,dm}{\int dm}.
\]
任何应用题都应做单位检查：$dA$、$ds$、$dV$、$dW$ 的量纲必须与最终量一致；若单位不对，说明微元翻译或变量选择发生了第一偏离。

\section{反常积分：先拆奇点，再谈无限对象}

反常积分不是“在定积分符号里直接代入无穷”。它是正常积分的极限，必须先识别所有奇点：无穷端点、有限瑕点、内部瑕点以及混合型奇点。

\subsection{定义与一票否决}

\[
\int_a^{+\infty}{f(x)}dx:=\lim_{R\to+\infty}\int_a^R{f(x)}dx,
\]
右端瑕点则定义为 $\lim_{t\to b^-}\int_a^t{f}$；内部瑕点 $c$ 必须拆成左右两个单侧积分。双无穷区间也要任取有限点 $c$ 分成两段，且两段分别收敛。

\begin{warnbox}{一票否决原则}
原反常积分收敛，当且仅当拆出的每个单奇点积分都收敛。不能把内部瑕点两侧的发散量相互抵消，也不能把普通收敛偷换成对称截断的 Cauchy 主值。
\end{warnbox}

\subsection{局部化}

远离奇点的闭区间上若函数连续，积分贡献是有限常数；敛散性只由各奇点的单侧邻域决定。因此可任意缩小有限瑕点附近的区间，或任意更换无穷积分的有限起点。这个局部化原则把复杂积分降为若干基准模型的比较。

\subsection{幂与幂对数母模型}

在 $0$ 附近，
\[
\int_0^1x^\alpha dx
\begin{cases}\text{收敛},&\alpha>-1,\\\text{发散},&\alpha\le-1,\end{cases}
\]
在 $+\infty$ 处，
\[
\int_1^{+\infty}x^\alpha dx
\begin{cases}\text{收敛},&\alpha<-1,\\\text{发散},&\alpha\ge-1.\end{cases}
\]
临界幂次用对数修正区分：
\[
\int_2^{+\infty}\frac{dx}{x{(\ln x)}^p}
\text{ 收敛当且仅当 }p>1.
\]
遇到更复杂的表达式，先提取主导幂次，再判断是否需要保留对数或振荡因子。

\section{判敛：最终保号与持续振荡}

\subsection{最终保号函数}

若在奇点邻域内 $0\le f\le g$ 且 $\int g$ 收敛，则 $\int f$ 收敛；若 $0\le g\le f$ 且 $\int g$ 发散，则 $\int f$ 发散。极限比较要求最终同号且
\[
\lim\frac{f(x)}{g(x)}=L\in(0,+\infty).
\]
渐近等价只能在同一局部、同一方向和同号条件下替换；若主部相消或函数持续变号，应改用更精细方法。

\subsection{振荡函数}

绝对收敛先检查 $\int|f|$。若绝对收敛失败但振荡抵消持续存在，可用 Dirichlet/Abel 型结构。Dirichlet 的典型接口是：振荡因子拥有一致有界的原函数，权重单调趋于零；Abel 则用一个已收敛的反常积分配合有界单调权重。分部积分把尾部压成边界项与可控积分。振荡并不自动收敛，必须写出“谁提供抵消、谁提供衰减”。

\subsection{Cauchy 准则与尾部}

无穷积分收敛等价于任意足够大的 $R<S$，尾部
\[
\int_R^S{f(x)}dx
\]
可以一致地小。这个准则是连续对象的“尾部稳定”版本，也是 H-B04 与数项级数建立接口的唯一安全入口。仅有 $f(x)\to0$ 不是一般的充分条件，甚至不是所有振荡积分的必要条件。

\section{反常积分的计算合法性}

计算反常积分时固定执行：
\begin{enumerate}
  \item 标出所有奇点和无穷方向；
  \item 按定义拆成单奇点段；
  \item 在有限截断区间上做换元、分部或原函数计算；
  \item 保留每个截断参数，分别取极限；
  \item 逐段检查极限是否有限，再汇总数值。
\end{enumerate}

\begin{warnbox}{不能先算发散项再相减}
若形式上出现 $\int f-\int g$，而两项分别发散，不能把它们先写成“$\infty-\infty$”后期待有限答案。必须先合并为同一个合法截断积分，或证明差函数本身在每个奇点附近可积；任何对称主值都必须明确标注其不是普通反常积分。
\end{warnbox}

\section{贯穿母例：同一旋转体的两种局部化}

把区域
\[
D=\{(x,y):0\le x\le1,\ 0\le y\le x(1-x)\}
\]
绕 $y$ 轴旋转。目标是体积，保留 $x$ 作坐标时，平行于旋转轴的竖直切片生成圆柱壳：半径为 $x$，壳高为 $x(1-x)$，故
\[
dV=2\pi x\cdot x(1-x)\,dx,
\qquad
V=2\pi\int_0^1x^2(1-x)\,dx=\frac{\pi}{6}.
\]
这完整运行了“对象与坐标 $\to$ 微元 $\to$ 局部贡献 $\to$ 区间 $\to$ 核验”：被积式量纲为长度平方，乘 $dx$ 后得到体积；结果为正且小于包围圆柱体积 $\pi/4$。

若改用垫圈法，就必须先承认 $y=x(1-x)$ 在 $[0,1]$ 上不是单射。对 $0\le y\le1/4$，外、内半径分别为
\[
R(y)=\frac{1+\sqrt{1-4y}}2,\qquad
r(y)=\frac{1-\sqrt{1-4y}}2,
\]
于是
\[
V=\pi\int_0^{1/4}\bigl(R^2-r^2\bigr)\,dy
=\pi\int_0^{1/4}\sqrt{1-4y}\,dy=\frac{\pi}{6}.
\]
两条独立路线相符。最典型的第一次偏离是只取一个反函数分支，把垫圈误画成圆盘；题面出现旋转体时，先比较切片方向与反解分支成本，再选择圆盘/垫圈或壳元。

\section{一元积分题的执行协议}

\subsection{有限区间协议}

\[
\begin{tikzpicture}[>=Stealth]
  \tikzset{
    flowstage/.style={draw=deep, fill=white, rounded corners=1.5mm, minimum height=0.72cm, align=center, font=\small\sffamily\bfseries\color{deep}, line width=0.6pt, inner xsep=6pt, inner ysep=3.5pt},
    flowarr/.style={->, >=Stealth, line width=0.75pt, draw=accent}
  }
  \node[flowstage] (n1) at (0,0) {1. 对象/定义域};
  \node[flowstage, right=0.4cm of n1] (n2) {2. 区间扫描};
  \node[flowstage, right=0.4cm of n2] (n3) {3. 对称与分段};
  \node[flowstage, below=0.6cm of n3] (n4) {4. 作差或换元};
  \node[flowstage, left=0.4cm of n4] (n5) {5. 计算};
  \node[flowstage, left=0.4cm of n5] (n6) {6. 符号/单位核验};

  \draw[flowarr] (n1) -- (n2);
  \draw[flowarr] (n2) -- (n3);
  \draw[flowarr] (n3) -- (n4);
  \draw[flowarr] (n4) -- (n5);
  \draw[flowarr] (n5) -- (n6);
\end{tikzpicture}
\]
若题目是比较或估值，计算步骤可以在作差、保序或平均高度处停止；若题目是应用建模，微元和单位检查优先于记忆公式。

\subsection{反常协议}

\[
\begin{tikzpicture}[>=Stealth]
  \tikzset{
    flowstage/.style={draw=deep, fill=white, rounded corners=1.5mm, minimum height=0.72cm, align=center, font=\small\sffamily\bfseries\color{deep}, line width=0.6pt, inner xsep=6pt, inner ysep=3.5pt},
    flowarr/.style={->, >=Stealth, line width=0.75pt, draw=accent}
  }
  \node[flowstage] (n1) at (0,0) {1. 奇点清单};
  \node[flowstage, right=0.4cm of n1] (n2) {2. 单点拆分};
  \node[flowstage, right=0.4cm of n2] (n3) {3. 局部化};
  \node[flowstage, below=0.6cm of n3] (n4) {4. 判敛或截断计算};
  \node[flowstage, left=0.4cm of n4] (n5) {5. 一票否决汇总};

  \draw[flowarr] (n1) -- (n2);
  \draw[flowarr] (n2) -- (n3);
  \draw[flowarr] (n3) -- (n4);
  \draw[flowarr] (n4) -- (n5);
\end{tikzpicture}
\]
得到一个有限数值前，不能省略“每段分别收敛”；得到一个发散结论时，只需找到一个发散段即可停止。

\subsection{结果验证}

至少选择一种独立检查：反向求导、端点代入、正负号、数量级、对称性、单位或数值近似。对反常积分，额外检查截断参数方向和是否误用了主值。

\section{高频失败路径：第一次偏离发生在哪里？}

\begin{enumerate}
  \item \kw{对象失败}：把原函数、定积分和变上限函数当成同一个对象，直接“积分与求导抵消”。
  \item \kw{区间失败}：忽略定义域断点、绝对值分段或换元后的上下限，导致积分对象已经改变。
  \item \kw{策略失败}：定积分先硬求原函数，错过奇偶、周期、翻转、保序和平均高度。
  \item \kw{微元失败}：切片方向、距离或密度翻译错误，公式看似熟悉但量纲不对。
  \item \kw{奇点失败}：没有列出全部瑕点，或把内部瑕点两侧的发散抵消成有限值。
  \item \kw{判敛失败}：只看 $f(x)\to0$，把最终保号比较法用于持续振荡，或把渐近等价用于主部相消。
  \item \kw{运算失败}：在反常积分中先写发散原函数差、先分部后取极限，遗漏边界项和参数方向。
  \item \kw{Owner 失败}：把变限积分可导性的细分类、级数判别或无限 Taylor 表示复制进本册。
\end{enumerate}

\section{传统章节与 Source Corpus 映射}

\begin{longtable}{L{2.55cm}L{5.35cm}L{\dimexpr\linewidth-7.90cm\relax}}
\toprule
\textbf{Source / 传统内容} & \textbf{进入本册} & \textbf{分流与拒绝复制}\\
\midrule
I-08 不定积分 & 原函数族、反导结构、凑微分/换元/分部/递推的决策机制 & 逐条公式表只保留母式；题目动作进入 Rules\\
I-08.1 三对象关系 & 原函数、Riemann 累积、变上限函数的对象分层 & Darboux、Thomae 和变限积分的细致正则性归 H-B03\\
I-09 定积分计算 & 线性、区间可加、对称、周期、换元、分部和 Newton\textendash{}Leibniz 接口 & 多重积分和完整参数积分另行建 Owner\\
I-09/I-10.1 比大小 & 保序、平均高度、作差/作商、局部放缩和结构扫描 & 高阶不等式的独立理论只作调用接口\\
I-10 微积分应用 & 微元法、面积、弧长、旋转体、功、压力、形心与单位核验 & 导数应用和多元应用不在本册展开\\
I-11 反常积分计算 & 奇点拆分、截断、合法换元/分部、发散项边界 & Gamma/Beta 仅作母例；特殊函数另行归属\\
I-12 反常积分判敛 & 局部化、幂/幂对数、比较、渐近、Dirichlet/Abel 和 Cauchy 尾部 & 连续/离散无限累积接口归 H-B04\\
II-09 无穷级数 & 只接收“有限截断—尾部稳定—无限对象”的接口 & 级数本体归 Topic10\\
\bottomrule
\end{longtable}

\section{一页压缩：忘掉公式后怎样重建？}

\begin{corebox}{一句话}
先问积分对象和合法区间，再把局部贡献变成有限累积；有限区间先做结构预处理，反常区间先拆奇点，最后才计算或取极限，并用独立核验确认结果。
\end{corebox}

\subsection{三条生成原则}
\begin{enumerate}
  \item \kw{对象原则}：原函数是反导，定积分是有向累积，变上限积分是新函数；名称相近不代表资格相同。
  \item \kw{有限优先原则}：任何无限对象都先截断成正常积分；任何计算变换都先在有限区间证明合法，再逐个取极限。
  \item \kw{密度原则}：比较积分、做物理建模或判断收敛时，真正要检查的是带微分因子的局部贡献及其符号/尺度。
\end{enumerate}

\subsection{重建问题}
\begin{enumerate}
  \item 当前讨论的是原函数、定积分、变上限函数，还是反常积分？
  \item 定义域、分段点、瑕点和无穷方向是否已经列全？
  \item 有限区间能否先用对称、周期、翻转、作差或保序降复杂度？
  \item 微元的局部贡献是什么，积分变量、区间和单位是否一致？
  \item 反常积分是否逐奇点拆分，是否存在一段发散而被错误抵消？
  \item 结果能否通过反向求导、端点、符号、数量级、单位或数值近似独立验证？
  \item 是否已经越界到 H-B03 的正则性、H-B04 的级数尾部或 H-B05 的无限 Taylor 表示？
\end{enumerate}

\begin{center}\rule{0.5\linewidth}{0.5pt}\end{center}

\appendix
\section{一元累积 Coverage 锚点}

\subsection{对象与资格速查}
\begin{center}
\begin{tabular}{L{2.55cm}L{5.15cm}L{\dimexpr\linewidth-7.70cm\relax}}
\toprule
对象 & 首要资格 & 输出\\
\midrule
原函数 & 在同一连通区间满足 $F'=f$ & 反导函数族 $F+C$\\
Riemann 定积分 & 被积函数在闭区间可积 & 有限有向累积\\
变上限积分 & $A(x)=\int_a^x{f}$，可导性另查正则性 & 新函数及其平均值商\\
反常积分 & 每个奇点/无穷方向的单侧极限分别存在 & 普通收敛、发散或明确主值\\
微元应用 & 局部贡献、切片方向、单位一致 & 面积、体积、功、矩等总量\\
\bottomrule
\end{tabular}
\end{center}

\subsection{母模型 Coverage}

本册覆盖：原函数与不定积分的对象边界，定积分性质与结构计算，积分比较与估值，微元建模，反常积分定义/拆分/计算/判敛。变限积分的精细可导性和原函数存在性反例由 H-B03 接收；连续无限累积与数项级数的判别翻译由 H-B04 接收；Taylor 级数由 Topic11/H-B05 接收。

\end{document}
